\documentclass[11pt,twoside]{amsart}
\usepackage{amssymb, amsmath, enumerate, palatino, hyperref}
\usepackage[normalem]{ulem}
\usepackage{fullpage}
\usepackage[T1]{fontenc}
\renewcommand{\labelitemi}{\guillemotright}
\usepackage{mathrsfs}

\usepackage{tikz}
\usepackage{tkz-graph}
\usetikzlibrary{positioning}
\tikzstyle{ball} = [circle,shading=ball, ball color=black,
    minimum size=1mm,inner sep=1.3pt]

\theoremstyle{plain}
\newtheorem{prop}{Proposition}%[section]
\newtheorem{lemma}[prop]{Lemma}
\newtheorem{thm}[prop]{Theorem}
\newtheorem{obs}[prop]{Observation}
\newtheorem{app}[prop]{Application}
\newtheorem*{MainThm}{Main Theorem}
\newtheorem{cor}[prop]{Corollary}
\newtheorem{conj}[prop]{Conjecture}
\theoremstyle{remark}
\newtheorem{rmk}[prop]{Remark}
\newtheorem{prob}{Problem}
\newtheorem{bonus}[prop]{Bonus Problem}
\newtheorem{exc}{Exercise}
\theoremstyle{definition}
\newtheorem{ex}[prop]{Example}
\theoremstyle{definition}
\newtheorem{defn}[prop]{Definition}

\newcommand{\RR}{\mathbb{R}}
\newcommand{\ZZ}{\mathbb{Z}}
\newcommand{\CC}{\mathbb{C}}
\newcommand{\NN}{\mathbb{N}}
\newcommand{\QQ}{\mathbb{Q}}
\newcommand{\PP}{\mathbb{P}}

\newcommand{\cC}{\mathscr{C}}
\newcommand{\Ob}{\operatorname{Ob}}
\newcommand{\Mor}{\operatorname{Mor}}

\newcommand{\Set}{\operatorname{Set}}
\newcommand{\FinSet}{\operatorname{FinSet}}
\newcommand{\Vect}{\operatorname{Vect}}
\newcommand{\FinVect}{\operatorname{FinVect}}
\newcommand{\Mat}{\operatorname{Mat}}
\newcommand{\Gp}{\operatorname{Gp}}
\newcommand{\FinGp}{\operatorname{FinGp}}
\newcommand{\AbGp}{\operatorname{AbGp}}
\newcommand{\Ring}{\operatorname{Ring}}
\newcommand{\CommRing}{\operatorname{CommRing}}
\newcommand{\Field}{\operatorname{Field}}
\newcommand{\Top}{\operatorname{Top}}
\newcommand{\Aut}{\operatorname{Aut}}
\newcommand{\id}{\operatorname{id}}

\newcommand{\defeq}{:=}

% For solutions.
\newenvironment{sol}{
  \medskip

  \noindent\textsc{Solution:}
}
{
  \medskip

}

\title{Math 113: Discrete Structures\\ Homework 17}
%\author{Your name} %Remove the leading % and put your name here if using this .tex file as a template for your homework.

\begin{document}
\maketitle

\noindent
\textbf{Due:} Monday, March 9 at 10pm.
\bigskip

A {\em spanning tree} of a connected graph~$G$ is a subgraph~$T$ such that~$T$ is
a tree and every vertex of~$G$ is on some edge of~$T$.  For instance,
if~$G$ is the triangle with vertices~$1,2,3$, then its spanning trees are:
\begin{center}
  \begin{tikzpicture}[scale=0.7]
    \node[ball, label={below left:$1$}] at (-1,0) (1) {};
    \node[ball, label={below right:$2$}] at (1,0) (2) {};
    \node[ball, label={above:$3$}] at (0,1.732) (3) {};
    \draw (1)--(3)--(2);
    \begin{scope}[xshift=4cm]
      \node[ball, label={below left:$1$}] at (-1,0) (1) {};
      \node[ball, label={below right:$2$}] at (1,0) (2) {};
      \node[ball, label={above:$3$}] at (0,1.732) (3) {};
      \draw (1)--(2)--(3);
    \end{scope}
    \begin{scope}[xshift=8cm]
      \node[ball, label={below left:$1$}] at (-1,0) (1) {};
      \node[ball, label={below right:$2$}] at (1,0) (2) {};
      \node[ball, label={above:$3$}] at (0,1.732) (3) {};
      \draw (2)--(1)--(3);
    \end{scope}
  \end{tikzpicture}
\end{center}
Recall that a {\em multigraph} is a graph in which multiple edges are
allowed.  For instance, the following graph has two edges connecting the
vertices~$1$ and~$3$:
\begin{center}
  \begin{tikzpicture}[scale=0.7]
    \node[ball, label={below left:$1$}] at (-1,0) (1) {};
    \node[ball, label={below right:$2$}] at (1,0) (2) {};
    \node[ball, label={above:$3$}] at (0,1.732) (3) {};
    \draw (1)--(2)--(3);
    \draw (1) edge[bend left] (3);
    \draw (1) edge[bend right] (3);
  \end{tikzpicture}
\end{center}
It has five spanning trees:
\begin{center}
  \begin{tikzpicture}[scale=0.7]
    \node[ball, label={below left:$1$}] at (-1,0) (1) {};
    \node[ball, label={below right:$2$}] at (1,0) (2) {};
    \node[ball, label={above:$3$}] at (0,1.732) (3) {};
    \draw (3)--(2);
    \draw (1) edge[bend left] (3);
    \begin{scope}[xshift=4cm]
      \node[ball, label={below left:$1$}] at (-1,0) (1) {};
      \node[ball, label={below right:$2$}] at (1,0) (2) {};
      \node[ball, label={above:$3$}] at (0,1.732) (3) {};
      \draw (3)--(2);
      \draw (1) edge[bend right] (3);
    \end{scope}
    \begin{scope}[xshift=8cm]
      \node[ball, label={below left:$1$}] at (-1,0) (1) {};
      \node[ball, label={below right:$2$}] at (1,0) (2) {};
      \node[ball, label={above:$3$}] at (0,1.732) (3) {};
      \draw (1)--(2)--(3);
    \end{scope}
    \begin{scope}[xshift=12cm]
      \node[ball, label={below left:$1$}] at (-1,0) (1) {};
      \node[ball, label={below right:$2$}] at (1,0) (2) {};
      \node[ball, label={above:$3$}] at (0,1.732) (3) {};
      \draw (2)--(1);
      \draw (1) edge[bend left] (3);
    \end{scope}
    \begin{scope}[xshift=16cm]
      \node[ball, label={below left:$1$}] at (-1,0) (1) {};
      \node[ball, label={below right:$2$}] at (1,0) (2) {};
      \node[ball, label={above:$3$}] at (0,1.732) (3) {};
      \draw (2)--(1);
      \draw (1) edge[bend right] (3);
    \end{scope}
  \end{tikzpicture}
\end{center}
\begin{prob}\label{prob: spanning trees}
  Draw all spanning trees of the following graph:
  \begin{center}
    \begin{tikzpicture}[scale=0.7]
      \node[ball, label={left:$1$}] at (-1,0) (1) {};
      \node[ball, label={right:$2$}] at (1,0) (2) {};
      \node[ball, label={above:$3$}] at (0,1.732) (3) {};
      \node[ball, label={below:$4$}] at (0,-1.732) (4) {};
      \draw (1)--(2)--(3)--(1)--(4)--(2);
    \end{tikzpicture}
  \end{center}
\end{prob}


\begin{prob}[deletion and contraction]\label{prob: deletion/contraction}
  Let~$G$ be a multigraph, and let~$e$ be an edge of~$G$.  Define~$G-e$ to be
  the graph obtained from~$G$ by removing the edge~$e$ (but retaining the
  endpoints of~$e$).  Let~$G/e$ be the graph obtained from~$G$ by
  ``contracting'' the edge~$e$.  To contract~$e$, remove~$e$ from~$G$ and then
  glue the endpoints of~$e$ together to make a single vertex from the two
  vertices.  If there were multiple edges between the endpoints of~$e$, loops
  will be formed, but for our purposes, we remove these loops as in the following:
  \begin{center}
    \begin{tikzpicture}[scale=0.7]
      \node[ball] at (-1,0) (1) {};
      \node[ball] at (1,0) (2) {};
      \node[ball] at (0,1.732) (3) {};
      %\node[ball, label={below left:$1$}] at (-1,0) (1) {};
      %\node[ball, label={below right:$2$}] at (1,0) (2) {};
      %\node[ball, label={above:$3$}] at (0,1.732) (3) {};
      \draw (1)--(2)--(3);
      \draw (1) edge[bend left] node[xshift=3,yshift=3,label={left:$e$}] {} (3);
      \draw (1) edge[bend right] (3);
      \draw[->] (2,1)--(3,1);
      \node at (0,-1.1) {$G$};
      \begin{scope}[xshift=6cm]
	\node[ball] at (-1,1) (1) {};
	\node[ball] at (1,1) (2) {};
	%\node[ball, label={below:$1$}] at (-1,1) (1) {};
	%\node[ball, label={below right:$2$}] at (1,1) (2) {};
        \draw (1) edge[bend right] (2);
        \draw (1) edge[bend left] (2);
	\draw (1) edge[loop above,looseness=30,in=120,out=60] node {} (1);
	\draw[->] (2.5,1)--(3.5,1);
      \end{scope}
      \begin{scope}[xshift=12cm]
	\node[ball] at (-1,1) (1) {};
	\node[ball] at (1,1) (2) {};
	%\node[ball, label={below:$1$}] at (-1,1) (1) {};
	%\node[ball, label={below right:$2$}] at (1,1) (2) {};
        \draw (1) edge[bend right] (2);
        \draw (1) edge[bend left] (2);
	\node at (0,-0.5) {$G/e$};
      \end{scope}
    \end{tikzpicture}
  \end{center}
  \begin{enumerate}[(a)]
    \item For an arbitrary connected multigraph~$G$, choose an edge~$e$ such
      that~$G-e$ is connected. Let~$T(G)$, $T(G-e)$, and~$T(G/e)$ denote the
      number of spanning trees of~$G$,~$G-e$, and~$G/e$, respectively.  The
      spanning trees of~$G$ come in two types: those that contain~$e$ and those
      that do not.  Use that idea to prove
      \[
       T(G) = T(G-e)+T(G/e).
      \]
    \item We can use the previous problem iteratively to count spanning trees.
      This is illustrated in the diagram below (at each stage, the edge chosen
	to delete and contract is dotted):
      \begin{center}
	\begin{tikzpicture}[scale=0.6]
	  \node[ball] at (-1,0) (1) {};
	  \node[ball] at (1,0) (2) {};
	  \node[ball] at (0,1.732) (3) {};
	  \draw (1)--(2)--(3);
	  \draw[very thick, dotted] (1) edge[bend left] (3);
	  \draw (1) edge[bend right] (3);
	  \draw (-2,-0.6) edge node[rectangle, fill=white, pos=0.5] {deletion}(-4.7,-2.9);
	  \draw (2,-0.6) edge node[rectangle, fill=white, pos=0.5] {contraction}(4.7,-2.9);
	  \begin{scope}[xshift=-6cm,yshift=-5cm]
	    \node[ball] at (-1,0) (1) {};
	    \node[ball] at (1,0) (2) {};
	    \node[ball] at (0,1.732) (3) {};
	    \draw (1)--(2)--(3);
	    \draw[very thick, dotted] (1) edge[bend right] (3);
	    \draw (-1.1,-0.6) edge node[rectangle, fill=white, pos=0.5] {deletion}(-3.4,-2.7);
	    \draw (1,-0.6) edge node[rectangle, fill=white, pos=0.5] {contraction}(2.5,-2.9);
	  \end{scope}
	  \begin{scope}[xshift=6cm,yshift=-5cm]
	    \node[ball] at (-1,0) (1) {};
	    \node[ball] at (1,0) (2) {};
	    \draw[very thick, dotted] (1) edge[bend left] (2);
	    \draw (1) edge[bend right] (2);
	    \draw (-1.1,-0.6) edge node[rectangle, fill=white, pos=0.5]
	    {deletion}(-2.8,-3.5);
	    \draw (1,-0.6) edge node[rectangle, fill=white, pos=0.5]
	    {contraction}(2.5,-3.5);
	  \end{scope}
	  \begin{scope}[xshift=-10cm,yshift=-10cm]
	    \node[ball] at (-1,0) (1) {};
	    \node[ball] at (1,0) (2) {};
	    \node[ball] at (0,1.732) (3) {};
	    \draw (1)--(2)--(3);
	  \end{scope}
	  \begin{scope}[xshift=-3cm,yshift=-10cm]
	    \node[ball] at (-1,0) (1) {};
	    \node[ball] at (1,0) (2) {};
	    \draw[very thick, dotted] (1) edge[bend left] (2);
	    \draw (1) edge[bend right] (2);
	    \draw (-0.8,-0.8) edge node[rectangle, fill=white, pos=0.5]
	    {deletion}(-2.2,-3.7);
	    \draw (1,-0.8) edge node[rectangle, fill=white, pos=0.5]
	    {contraction}(2.5,-3.7);
	  \end{scope}
	  \begin{scope}[xshift=3cm,yshift=-10cm]
	    \node[ball] at (-1,0) (1) {};
	    \node[ball] at (1,0) (2) {};
	    \draw (1) edge[bend right] (2);
	  \end{scope}
	  \begin{scope}[xshift=10cm,yshift=-10cm]
	    \node[ball] at (-1,0) (1) {};
	  \end{scope}
	  \begin{scope}[xshift=-5.4cm,yshift=-14.8cm]
	    \node[ball] at (-1,0) (1) {};
	    \node[ball] at (1,0) (2) {};
	    \draw (1) edge[bend right] (2);
	  \end{scope}
	  \begin{scope}[xshift=1cm,yshift=-14.8cm]
	    \node[ball] at (-1,0) (1) {};
	  \end{scope}
	\end{tikzpicture}
      \end{center}
  \end{enumerate}
  \medskip

  We stop in this deletion-contraction process when there are no edges left
  whose removal would leave a connected graph.
  Along the bottom, there are~$5$ trees (a single isolated vertex is considered
  to be a tree, too).  The previous part of this problem implies there
  are~$5$ spanning trees of the original graph.  These are the~$5$ spanning
  trees we saw earlier.  
  
  Make a similar diagram for the graph in
  Problem~\ref{prob: spanning trees}.  (This diagram should verify the number of
  spanning trees you found earlier.)\footnote{The first part of
    Problem~\ref{prob: deletion/contraction} implies the amazing fact that
    number of trees at the bottom of the diagram is independent of the choices
    of edges made in constructing the diagram!}
\end{prob}
\end{document}
 
